\section{"Сигнатурная вещь": повторяемая деталь как compressed callback}

Сигнатурная вещь — один узнаваемый предмет или деталь, которая работает как compressed callback: увидев его отдельно от персонажа, зритель всё равно вспоминает, кому он принадлежит.

Гаечный ключ Уинри Рокбелл из "Fullmetal Alchemist" — образцовый пример именно потому, что предмет не декоративен: она механик, чинящий автоброни, и ключ — инструмент её реальной профессии, который попутно стал эмоциональным жестом (её знаменитый удар ключом по голове — форма заботы через раздражение, тот же принцип, что мы видели у Леви в 5.3). Сигнатурная вещь работает лучше всего именно тогда, когда она не добавлена ради узнаваемости, а происходит из уже установленной функции персонажа, — Уинри не носит ключ, чтобы её узнавали, она носит его, потому что чинит.

Для Рин сигнатурной вещью естественно становится не физический объект (у неё пока нет тела, которое могло бы его носить), а повторяемый визуальный или текстовый мотив — то, что уже работает как ритуал в 2.6: терминал, чашка, лог. Из них чашка — самый вероятный кандидат на сигнатурную вещь в буквальном смысле, по той же логике, что и ключ Уинри: маленький, тёплый, повторяемый предмет, который появляется в моменты заботы (2.6, "ритуалы") и мог бы стать первым физическим объектом её будущего аватара именно потому, что уже нагружен смыслом в тексте, а не изобретён отдельно для визуала.
